\section{ETL de CSV a Supabase}

Es el pipeline histórico: convierte la carpeta de CSV crudos en dos tablas normalizadas.
No es un script de una sola vez, sino un proceso reproducible que se vuelve a correr
cada vez que aparecen archivos nuevos.

\begin{figure}[htbp]
\centering
\resizebox{\textwidth}{!}{% Figura: pipeline ETL de CSV, etapa por etapa, con el modulo responsable.
\begin{tikzpicture}[node distance=9mm]

\node[fuente,etapa] (a) {CSV crudo\\\nota{ragged, 13 esquemas}};
\node[etapa,right=of a] (b) {\textbf{extract}\\\nota{leer + tipar}\\\nota{+ normalizar cabeceras}};
\node[etapa,right=of b] (c) {\textbf{transform}\\\nota{separar flujos}\\\nota{+ remuestrear}};
\node[etapa,right=of c] (d) {\textbf{load}\\\nota{UPSERT por}\\\nota{\cd{timestamp}}};

\node[almacen,font=\scriptsize,text width=38mm,right=12mm of d] (t1) {\cd{monitoreo\_sc\_electrico}\\\nota{5 min}};
\node[almacen,font=\scriptsize,text width=38mm,below=8mm of t1] (t2) {\cd{radiacion\_sc\_15s}\\\nota{15 s}};

\draw[flecha] (a) -- (b);
\draw[flecha] (b) -- node[etiq,above]{\cd{slugify}} (c);
\draw[flecha] (c) -- (d);
\draw[flecha] (d.east) -- ++(4mm,0) |- (t1);
\draw[flecha] (d.east) -- ++(4mm,0) |- (t2);

% Estado / idempotencia
\node[comp,text width=30mm,below=14mm of b] (log) {\cd{\_ingest\_log}\\\nota{md5 por archivo}};
\draw[flechad] (a) |- node[etiq,pos=0.72,below]{¿cambió?} (log);
\draw[flechad] (log) -| (d.south);

% Capa solar
\node[comp,text width=36mm,below=14mm of t2] (solar) {\textbf{calibracion} + \textbf{performance}\\\nota{pvlib: clear-sky y POA}};
\draw[flecha] (t2) -- (solar);
\node[almacen,font=\scriptsize,text width=36mm,left=12mm of solar] (t3) {\cd{radiacion\_sc\_clearsky}\\ \cd{radiacion\_sc\_poa}};
\draw[flecha] (solar) -- (t3);

\end{tikzpicture}
}
\caption{Pipeline de CSV. El \cd{\_ingest\_log} decide qué archivos ni siquiera se abren;
la capa solar (pvlib) se refresca al final de cada corrida.}
\end{figure}

\subsection{Una sola lista, y todo lo demás se deduce de ella}

El problema de fondo de estos CSV es que el mismo dato viene escrito de muchas maneras
distintas. El voltaje del primer arreglo aparece como \cd{vpv1} en unos archivos, como
\cd{Voltaje PV1 [V]} en otros y como \cd{voltaje\_pv1\_v} en los más nuevos, y a eso hay
que sumarle las erratas y los acentos que se fueron colando. En total, unas setenta formas
distintas de nombrar apenas una docena de cosas.

La salida fácil sería escribir a mano la lista de nombres buenos y repetirla en cada lugar
donde hace falta: al crear las tablas, al insertar, al agrupar por intervalos de tiempo.
Es también la forma más segura de que dentro de seis meses alguien agregue una variable en
tres de esos cuatro lugares, se olvide del cuarto, y nadie se entere hasta que los números
no cuadren.

Por eso acá hay \textbf{un solo lugar} donde se declara qué se mide: un diccionario con
una entrada por cada cosa medida, que dice cuál es su nombre bueno. Todo lo demás sale de
ahí, incluido el SQL que crea la base de datos, que es un archivo \textbf{generado} y no
se escribe a mano.

Para quien tenga que trabajar con esto mañana, la diferencia se nota en dos cosas:

\begin{itemize}[leftmargin=*,itemsep=2pt]
  \item \textbf{Agregar una variable es escribir un renglón.} Aparece sola en el esquema,
    en la tabla y al agrupar por intervalos. No hay que acordarse de nada más.
  \item \textbf{Un nombre nuevo para algo que ya existe no da trabajo.} Si el archivo del
    mes que viene escribe la misma medición con otra errata, otro acento u otras unidades,
    el sistema la reconoce y la manda al mismo sitio, sin tocar una línea de código. Así
    fue como esas setenta variantes quedaron reducidas a una lista limpia.
\end{itemize}

En el código, el diccionario es \cd{CONCEPT\_MAP} y vive en el módulo \cd{normalize}.

\subsection{Módulos}

\begin{center}
\small
\begin{tabular}{@{}p{26mm}p{114mm}@{}}
\toprule
\textbf{Módulo} & \textbf{Responsabilidad única} \\
\midrule
\cd{normalize} & \cd{slugify} normaliza acentos, unidades y mayúsculas; \cd{CONCEPT\_MAP}
  traduce el slug al nombre canónico. \\
\cd{schemas} & Deriva \cd{CANONICAL\_COLUMNS}, infiere etiquetas del nombre y decide el
  método de agregación por etiqueta. \\
\cd{extract} & Lee el CSV tolerando filas irregulares, normaliza cabeceras, tipa y parsea
  el \cd{timestamp}. \\
\cd{transform} & \cd{split\_streams} separa el flujo eléctrico del de radiación
  \emph{por columnas}; remuestrea a 5\,min y 15\,s. \textbf{Sin limpieza}. \\
\cd{ddl} & Genera tablas, vistas de corrección y calibración, diccionario y sentencias RLS. \\
\cd{clearsky} & GHI de cielo despejado con pvlib (modelo Ineichen). \\
\cd{calibracion} & Puebla \cd{radiacion\_sc\_clearsky} de forma incremental. \\
\cd{performance} & Calcula POA por arreglo (transposición Erbs + isotrópica, modelo
  bifacial de dos planos) y puebla \cd{radiacion\_sc\_poa}. \\
\cd{load} & UPSERT por \cd{timestamp} en cada tabla. \\
\cd{state} & Registro md5 en \cd{\_ingest\_log}. \\
\cd{pipeline} & Orquesta \emph{extract $\to$ transform $\to$ load} y el refresco solar. \\
\cd{cli} & Menú interactivo numerado. \\
\bottomrule
\end{tabular}
\end{center}

\subsection{Idempotencia}

Dos niveles, deliberadamente redundantes:

\begin{itemize}[leftmargin=*,itemsep=2pt]
  \item \textbf{Nivel archivo}: \cd{\_ingest\_log} guarda el md5 de cada CSV. Si no
    cambió, ni se abre.
  \item \textbf{Nivel fila}: \cd{timestamp} es PRIMARY KEY con
    \cd{ON CONFLICT DO UPDATE}. Reprocesar nunca duplica, aunque el md5 falle.
\end{itemize}

Un archivo que revienta no frena la corrida: se registra el fallo, se hace \emph{rollback}
de esa transacción y el pipeline sigue con el siguiente. La capa solar va en su propio
\cd{try}, para que un problema de pvlib no invalide una carga de datos que ya funcionó.

\subsection{Uso}

Un solo comando, menú interactivo:

\begin{center}\cd{python3 main.py}\end{center}

Las opciones van en orden de confianza creciente: auditar el dataset, ejecutar en seco
sin tocar la base, generar el DDL, aplicarlo en Supabase, y recién ahí cargar datos
(incremental o reprocesando todo). Para archivos nuevos basta soltarlos en
\cd{dataset/Monitoreo-AgroVoltaic-SC-NEW/} y repetir la carga incremental.

\begin{aviso}[title=Lo que este ETL todavía no hace]
La \textbf{separación fina de filas mezcladas} (el llamado Paso 2) sigue pendiente. Hoy
se conservan las filas que coinciden con la cabecera del archivo y las irregulares se
saltan con registro en el log. El criterio acordado es recuperar lo que se pueda y
aceptar huecos; existe un par de archivos original/corregido como referencia.
\end{aviso}
